\subsection*{Chapter 14: Evaluating LLMs and Reasoning Models}

\paragraph{Exercise 14.1} Typical findings: (i) \emph{false negative} for unordered answers
(``2, 3'' vs.\ ``3, 2'' for ``all solutions''), because tuples are compared in order; fixing it
globally would break ordered pairs, so make it opt-in per problem (Exercise 14.2). (ii) \emph{false
negative} for units the reference keeps (``5\,\textbackslash text\{ cm\}'' vs.\ ``5'' unwraps the
\texttt{\textbackslash text}, leaving ``5 cm'', which SymPy parses as $5\,c\,m$); stripping known
unit words fixes it but risks deleting variables named $m$. (iii) \emph{false positive} from the
assignment rule: ``$y = 5$'' accepted for a question about $x$; restricting the rule to the variable
named in the problem fixes it at the cost of parsing the problem. Every fix is a policy change ---
add a regression test for each.

\paragraph{Exercise 14.2} Normalize both sides, split into components, then match as a multiset:
for each gold component find an unused predicted component with \texttt{equality\_check}. Naively
that is $O(k^2)$ SymPy calls. Cheaper: first bucket components by a canonical numeric key (evaluate
each with \texttt{sympy.N} to, say, 12 significant digits, or use the normalized string), match
identical keys in $O(k \log k)$, and fall back to pairwise symbolic checks only for the
unmatched remainder.

\paragraph{Exercise 14.3} Record \texttt{hit\_eos} (generation ended with EOS before the budget) and
\texttt{has\_box}. Then: \texttt{has\_box} and wrong $\Rightarrow$ the model is wrong;
\texttt{not hit\_eos} and \texttt{not has\_box} $\Rightarrow$ the budget truncated the chain of
thought (raise \texttt{max\_new\_tokens}); \texttt{hit\_eos} and \texttt{not has\_box}
$\Rightarrow$ a format failure (adjust the prompt or extraction fallback). Reporting the three rates
alongside accuracy explains most surprising accuracy changes.

\paragraph{Exercise 14.4} (a) A: 205/500, Wilson 95\% CI $[0.368, 0.454]$; B: 223/500, CI
$[0.403, 0.490]$. The intervals overlap, so independent intervals suggest no clear difference.
(b) McNemar uses only discordant items: $b = 30$, $c = 48$. With continuity correction
$\chi^2 = (|b-c|-1)^2/(b+c) = 289/78 = 3.71$, $p \approx 0.054$; without correction $\chi^2 = 4.15$,
$p \approx 0.042$; the exact binomial test gives $p \approx 0.054$. The paired test is the right one
and is far more sensitive than comparing intervals, because per-item difficulty cancels; here the
evidence is borderline, and a second seed or more items would be needed before claiming B is better.

\paragraph{Exercise 14.5} Useful row-steps are $\sum_i \ell_i$ out of $B \max_i \ell_i$, so the
wasted fraction is $1 - \bar\ell / \max_i \ell_i$. For i.i.d.\ exponential lengths with mean
$\mu = 600$, $\E[\max] = \mu H_B$ with $H_{64} \approx 4.744$, i.e.\ $\approx 2846$ tokens, so
waste $\approx 1 - 600/2846 \approx 79\%$. Even with a budget cap that truncates the tail, static
batching of long reasoning traces wastes most of the batch --- the motivation for continuous
batching (\cref{ch:serving}).

\paragraph{Exercise 14.6} For GPT-2 and Qwen BPE, \texttt{"A"} and \texttt{" A"} are distinct ids,
and \texttt{encode("Answer: A")[-1]} is the space-prefixed variant. Summing the softmax probability of
both variants per letter (computed in float32) is the robust choice; on base models the
single-variant (no-space) version can disagree on a noticeable fraction of questions because almost
all of the mass sits on the space-prefixed token.

\paragraph{Exercise 14.7} Sample votes from fixed strengths (e.g.\ Elo-scale 1000, 1050, 1100, 1200)
by drawing random pairs and a winner with probability $p_i/(p_i+p_j)$. Online Elo with $K = 32$
produces visibly different ratings for different vote orders (spreads of tens of points), most of it
driven by the last few hundred votes; BT is order-invariant by construction, and its bootstrap
interval reflects only sampling noise. Lowering $K$ reduces Elo's order dependence but slows
convergence.

\paragraph{Exercise 14.8} Write $\ell(p) = \sum_i W_i \log p_i - \sum_{i<j} n_{ij}\log(p_i+p_j)$.
Since $-\log x \ge 1 - \log y - x/y$ for $x, y > 0$ (concavity of $\log$), at the current iterate
$p^{(k)}$ the function
$Q(p) = \sum_i W_i\log p_i - \sum_{i<j} n_{ij}\bigl(\log(p^{(k)}_i+p^{(k)}_j) + \frac{p_i+p_j}{p^{(k)}_i+p^{(k)}_j} - 1\bigr)$
minorizes $\ell$ and touches it at $p^{(k)}$. $Q$ separates over $i$; setting
$\partial Q/\partial p_i = W_i/p_i - \sum_{j\ne i} n_{ij}/(p^{(k)}_i+p^{(k)}_j) = 0$ gives exactly the
MM update. Hence $\ell(p^{(k+1)}) \ge Q(p^{(k+1)}) \ge Q(p^{(k)}) = \ell(p^{(k)})$. Rescaling
does not change $\ell$, which is invariant to multiplying all $p_i$ by a constant.

\paragraph{Exercise 14.9} Position bias appears as verdicts that flip when the order of the two
responses is swapped; with small local judges flip rates of tens of percent are common. Debias by
evaluating both orders and (a) counting a win only if it holds in both (ties otherwise), or (b)
averaging the two scores/probabilities. Report the flip rate as a judge-quality metric.

\paragraph{Exercise 14.10} Work: $20 \times 500 \times 8 = \num{80000}$ generations; at an assumed
mean of 2000 tokens that is $1.6\times10^8$ decode tokens. A 7B model is 14\,GB in bf16 and fits on
one GPU, so use data parallelism (one replica per GPU; tensor parallelism only adds communication).
KV for a Llama-style 7B with GQA (32 layers, 8 KV heads, head dim 128) is $\approx 131$\,KB/token; at
batch 128 and 2000-token average context that is $\approx 34$\,GB. Each step moves
$\approx 14 + 34 = 48$\,GB, i.e.\ $\approx 14$\,ms on an H100, producing 128 tokens: $\approx 9$k
tokens/s per GPU, $\approx 72$k tokens/s for 8 GPUs, so $\approx 37$ minutes at the roofline; with
realistic 2--3$\times$ overheads, 1.5--2 hours. That fits in 6 hours. Grading 80k answers at
$\sim$10\,ms per SymPy check is $\sim$800 CPU-seconds --- a few cores keep up easily, provided
grading runs concurrently with generation.
